\begin{figure}[h]
    \centering
    \begin{tikzpicture}
        \begin{axis}[
            width=8.2cm, height=7cm,
            xlabel={$10^{-3} J \, [kg \, m^2]$},
            ylabel={Gęstość prawdopodobieństwa},
            grid=major,
            scaled x ticks=false,
            x tick label style={/pgf/number format/fixed, /pgf/number format/precision=4},
            legend style={at={(0.02,0.98)}, anchor=north west}
        ]

        % --- Metoda I ---
        \def\muOne{9.59}
        \def\sigmaOne{0.25}

        % --- Metoda II ---
        \def\muTwo{8.3}
        \def\sigmaTwo{4.2}

        \def\medium{(2 * \muOne + 1 * \muTwo) / 3}
        \def\range{4 * \sigmaOne}

        % --- Metoda I ---
        \addplot [
            domain=\medium - \range : \medium + \range,
            samples=1000,
            color=red,
            thick,
            smooth
        ]
        {1/(\sigmaOne*sqrt(2*pi)) * exp(-((x-\muOne)^2)/(2*\sigmaOne^2))};
        \addlegendentry{Metoda I}

        % --- Metoda II ---
        \addplot [
            domain=\medium - \range : \medium + \range,
            samples=1000,
            color=blue,
            thick,
            smooth
        ]
        {1/(\sigmaTwo*sqrt(2*pi)) * exp(-((x-\muTwo)^2)/(2*\sigmaTwo^2))};
        \addlegendentry{Metoda II}

        % --- Punkty --- %
        \pgfmathsetmacro{\peakOne}{1/(\sigmaOne*sqrt(2*pi))}
        \draw [dashed, black!70] (axis cs:\muOne, 0) -- (axis cs:\muOne, \peakOne);
        \addplot[mark=*, mark size=1.5pt, red] coordinates {(\muOne, \peakOne)};
        \addplot[mark=*, mark size=1.5pt, red] coordinates {(\muOne, 0)};
        
        \pgfmathsetmacro{\peakTwo}{1/(\sigmaTwo*sqrt(2*pi))}
        \draw [dashed, black!70] (axis cs:\muTwo, 0) -- (axis cs:\muTwo, \peakTwo);
        \addplot[mark=*, mark size=1.5pt, blue] coordinates {(\muTwo, \peakTwo)};
        \addplot[mark=*, mark size=1.5pt, blue] coordinates {(\muTwo, 0)};

        \end{axis}
    \end{tikzpicture}
    \caption{Porównanie gęstości prawdopodobieństwa uzyskanych momentów bezwładności dla obu metod wraz z zaznaczonymi wartościami oczekiwanymi (wąska domena).}
    \label{img:gestosc_prawdop_porownanie_wasko}
\end{figure}
